\begin{summarybox}
\textbf{一句话核心：沿抛物线的函数方向消元，先判断实点，再用韦达。}

设 $p>0,s=\pm1$，直线 $Ax+By+C=0$，$A^2+B^2>0$。

\textbf{左右开口，优先看纵坐标：}
\[
 y^2=2spx\ \Longrightarrow\ Ay^2+2spBy+2spC=0.
\]
\[
 y_1+y_2=-\frac{2spB}{A},\qquad y_1y_2=\frac{2spC}{A},\qquad
 H=p^2B^2-2spAC\quad(A\ne0).
\]
\textbf{上下开口，优先看横坐标：}
\[
 x^2=2spy\ \Longrightarrow\ Bx^2+2spAx+2spC=0.
\]
\[
 x_1+x_2=-\frac{2spA}{B},\qquad x_1x_2=\frac{2spC}{B},\qquad
 K=p^2A^2-2spBC\quad(B\ne0).
\]
\textbf{判断顺序：}
\begin{enumerate}[leftmargin=*,itemsep=.4em]
 \item 左右开口若 $A=0$，或上下开口若 $B=0$，均为轴向直线：恰一个交点，非相切。
 \item 否则查看 $H$ 或 $K$：正值两点，零值相切，负值无实交点。
 \item 确有两点后，坐标和的一半给出中点；根差配合勾股关系给出弦长。
\end{enumerate}
\textbf{横坐标固定形式仍可速查。} 左右开口时
\[
 A^2x^2+(2AC-2spB^2)x+C^2=0.
\]
$A\ne0$ 才用二次韦达；$B=0$ 时它只记录固定横坐标，不能单独证明实交点存在。

\textbf{最后四查：} 认准 $2p$ 与开口符号；不除以零；平方后回原直线；参数题舍去相切候选。
\end{summarybox}
